% Vietnamese translation for OI Public Library
% Source-File: IOI中国国家候选队论文/2003/方奇/方奇.ppt
\documentclass[11pt]{article}
\usepackage{oipl-vn}

\title{Bản trình chiếu: Tô màu và cấu tạo trên bàn cờ}
\author{Fang Qi}
\date{2003}

\begin{document}
\maketitle

\origtitle{Slide companion for coloring and construction on chessboards}
\sourcepath{IOI China candidate team papers / 2003 / Fang Qi / Fang Qi.ppt}

\section{Phạm vi bản dịch}

Bản trình chiếu đi kèm bài viết về phương pháp tô màu và cấu tạo. Phần chữ
khôi phục được chứa các ký hiệu \(m,n,k\), ô \(a(i,j)\), các đại lượng
\(T_{\text{free}},T_{\text{black}},T_{\text{white}}\), và nhiều sơ đồ bàn cờ.

\section{Phương pháp tô màu}

Với bàn cờ \(m\times n\), mỗi ô được ký hiệu \(a(i,j)\). Ta tô màu theo tính
chẵn lẻ của \(i+j\), hoặc theo một chu kỳ khác phù hợp với quân cờ hay mảnh
ghép đang xét. Khi một cấu hình chiếm các ô theo mẫu cố định, so sánh số ô
đen và trắng cho ta điều kiện cần.

\section{Phương pháp cấu tạo}

Sau khi có điều kiện cần, các slide chuyển sang xây dựng cấu hình đạt được.
Những đoạn như
\[
  m=m_1k+r,\quad 0<r<k,\qquad n=n_1k+s,\quad 0<s<k
\]
cho thấy bài toán được tách theo phần nguyên và phần dư của kích thước bàn
cờ. Các sơ đồ kiểu \(A-C-\), \(-D-B\) biểu diễn mẫu lát hoặc mẫu điền lặp lại.

\section{Liên hệ với bài viết}

Bản bài viết đã dịch giải thích đầy đủ hơn các bài con. Bản trình chiếu chủ
yếu phục vụ minh họa trực quan cho hai bước: dùng tô màu để chứng minh ràng
buộc, rồi dùng cấu tạo tuần hoàn để đạt ràng buộc đó.

\end{document}
